% ===========================================================================
%  sviluppatore/02-ambiente/02-avvio.tex
% ===========================================================================
\section{I quattro modi di avviarla}\label{sez:avvio}

Nella radice del progetto ci sono quattro punti d'ingresso, e fanno cose
diverse. Si lanciano tutti \emph{dalla radice}, mai da dentro una
sottocartella: gli import sono assoluti rispetto a quella.

\begin{center}
\begin{tabular}{@{}L{52mm}L{18mm}L{66mm}@{}}
\toprule
\textbf{Comando} & \textbf{Porta} & \textbf{Che cosa avvia} \\
\midrule
\rigaTenue
\file{python scripts/run\_dashboard.py}
  & \num{5000} & L'interfaccia completa. È questo che si usa nel
                 \num{99}\,\% dei casi. \\
\file{python run\_api.py}
  & \num{8000} & Solo l'API FastAPI del motore. Nessuna pagina, solo JSON. \\
\rigaTenue
\file{python desktop\_app.py}
  & \num{5000} & La stessa dashboard, in una finestra nativa invece che in una
                 scheda del browser. Richiede \pkg{pywebview}. \\
\file{npm run dev} (in \file{sito/})
  & \num{3000} & Il sito. Indipendente dai tre precedenti. \\
\bottomrule
\end{tabular}
\end{center}

\subsection{La dashboard}

\begin{codiceInterno}{console}{Avvio normale}{}
python scripts/run_dashboard.py
# http://127.0.0.1:5000
\end{codiceInterno}

\noindent
Host e porta si cambiano con le variabili d'ambiente \cod{FLASK\_HOST} e
\cod{FLASK\_PORT}, oppure nei campi corrispondenti della configurazione. Il
default è \cod{127.0.0.1}, cioè raggiungibile solo dalla macchina stessa, e
questo non è un dettaglio di comodità.

\begin{attenzione}{La dashboard non va esposta in rete}
Quasi nessuna rotta della dashboard richiede un login, ed è una scelta
deliberata: l'applicazione gira sul computer di chi la usa, e chiedere un
account per entrare in casa propria non avrebbe senso. Il prezzo è che chi
raggiunge la porta \num{5000} può fare tutto: lanciare backtest, leggere le
credenziali del broker, scrivere file. Mettere \cod{FLASK\_HOST=0.0.0.0} su
una macchina raggiungibile da fuori significa consegnarla. Le uniche aree
protette sono quelle che toccano il sito --- marketplace, profilo,
abbonamento, amministrazione --- e il \vediSez{permessi} spiega
esattamente quali e perché.
\end{attenzione}

\subsection{L'API del motore}

\begin{codiceInterno}{console}{Con ricaricamento automatico}{}
uvicorn api.app:app --reload
# http://127.0.0.1:8000/docs  -> la documentazione interattiva generata
\end{codiceInterno}

\noindent
Gira \emph{accanto} alla dashboard, non al suo posto. Per lavorare
sull'interfaccia non serve; serve quando si tocca un caso d'uso, perché
l'API lo espone senza passare da Flask ed è il modo più rapido per provarlo.

\subsection{Il sito}

\begin{codiceInterno}{console}{Dalla cartella del sito}{}
cd sito
npm install
npm run dev
# http://localhost:3000
\end{codiceInterno}

\noindent
Il sito ha bisogno di un Postgres e delle proprie variabili d'ambiente
(\cod{DATABASE\_URL}, \cod{NEXTAUTH\_SECRET}, le chiavi Stripe e SMTP). In
produzione queste non stanno in un file: si gestiscono dal pannello del
sito, ed è lo stesso meccanismo che poi le passa alla metà locale attraverso
il canale~\num{2} (\vediSez{due-meta}).

\begin{pratica}{Il giro minimo per vedere qualcosa}
\begin{enumerate}
  \item \file{pip install -r requirements.txt}
  \item \file{python scripts/download\_d1.py} --- scarica un po' di dati
        giornalieri (\vediSez{dati}).
  \item \file{python scripts/run\_dashboard.py}
  \item Nel browser: \voce{Run} \freccia{} una strategia \freccia{}
        \pulsante{Run backtest}. Alla fine compare uno stream in
        \voce{Streams}.
\end{enumerate}
Senza il passo~\num{2} l'applicazione parte e le pagine sono vuote: non è un
errore, è che non c'è niente da mostrare.
\end{pratica}
